%% Wiki: [[2026.01.29 - Mathematics: Optimal Poset-Feasible Padding: Knapsack, Transport, and Dual Certificates for TV Privacy]]

\documentclass[11pt]{article}
\usepackage[margin=1in]{geometry}
\usepackage{amsmath,amssymb,amsfonts,amsthm}
\usepackage{microtype}
\usepackage{hyperref}
\usepackage{enumitem}
\usepackage{framed}

\hypersetup{colorlinks=true,linkcolor=blue,citecolor=blue,urlcolor=blue}

\newtheorem{theorem}{Theorem}
\newtheorem{lemma}[theorem]{Lemma}
\theoremstyle{definition}
\newtheorem{definition}[theorem]{Definition}
\newtheorem{example}[theorem]{Example}
\theoremstyle{remark}
\newtheorem{remark}[theorem]{Remark}

\newcommand{\TV}{\mathrm{TV}}
\newcommand{\E}{\mathbb{E}}
\newcommand{\Prb}{\mathbb{P}}


\title{Optimal Poset-Feasible Padding: Knapsack, Transport, and Dual Certificates for TV Privacy\\
\large Lean addendum cut (assumes prior papers open)}
\author{}
\date{}

\begin{document}
\maketitle

\noindent\textbf{Canonical wiki page:} \texttt{[[2026.01.29 - Mathematics: Optimal Poset-Feasible Padding: Knapsack, Transport, and Dual Certificates for TV Privacy]]}\\


\vspace{-0.75em}
\begin{abstract}
oindent
\textbf{Assumption: non-self-contained.} Readers have immediate access to:
\emph{Spectral Anonymity and Optimal Distinguishing Bounds for Random-Walk Delegation in (Possibly Directed) Overlays},
\emph{Prefix Capacities and Separation Cuts for Shared-Kernel Termination-Time Hiding},
\emph{stop\_time\_unique\_addendum\_v3},
and \emph{many\_testing\_pub\_tight\_section2and4}.
This addendum records only the parts that are genuinely \emph{poset-native} and that we use later:
\emph{(i)} the condition under which top-censoring collapses to fractional knapsack,
\emph{(ii)} a pooling curve for two-point swap worlds together with global join-optimality on join-semilattices,
\emph{(iii)} at perfect transcript privacy ($\varepsilon=0$), a join-transport lower bound and a sharp likelihood-ratio implementability criterion,
and \emph{(iv)} for general $\varepsilon$, a compact dual exposing slopes/kinks plus small-menu/sparsity consequences.
All motivation, TV-as-testing interpretation, chain/stop-time special cases, and certificate machinery are deferred to the four prior papers (titles above).
\end{abstract}

\paragraph{Minimal model.}
Let $(Y,\preceq)$ be a finite poset of transcripts and $X\subseteq Y$ intrinsic states.
A kernel $K(y\mid x)$ is \emph{poset-feasible} if $K(y\mid x)=0$ unless $x\preceq y$ and $\sum_{y\succeq x}K(y\mid x)=1$.
Two worlds $b\in\{0,1\}$ have priors $\pi_b$ on $X$; $K$ induces $P_b(y)=\sum_x \pi_b(x)K(y\mid x)$.
We minimize expected monotone cost subject to $\TV(P_0,P_1)\le \varepsilon$ (and feature-family variants).

\vspace{0.25em}\noindent
\textbf{What is new here (one-screen summary).}
\begin{itemize}[leftmargin=*]
\item \textbf{Poset abstraction:} $x\preceq y$ formalizes add-only transcripts beyond total orders; earlier papers cover the chain/stop-time case.
\item \textbf{Non-aligned structure:} optimal kernels can \emph{pool} at intermediate upper bounds; for swap worlds on join-semilattices, the join is globally optimal.
\item \textbf{Perfect privacy:} a two-kernel coupling relaxation equals \emph{join-transport}; implementability by a single kernel is exactly likelihood-ratio invariance.
\item \textbf{Dual certificates at general $\varepsilon$:} a compact dual yields slope/kink certificates and small-menu/sparsity bounds.
\end{itemize}

\noindent
\textbf{Notation bridge.} Let $\bar\pi=\tfrac12(\pi_0+\pi_1)$ and $w(x):=\pi_0(x)-\pi_1(x)$.
In a few imported theorem statements below, the same object may appear as $\Delta(x)$ or $d(x)$; read all of them as $w(x)$.
For swap worlds we reserve $\delta$ for the \emph{scalar} imbalance between the two swapped masses.
We write $c$ for a monotone cost (if $u\preceq v$ then $c(u)\le c(v)$).

\vspace{0.25em}\noindent
\textbf{Standing conditions (only when invoked).}
Unless stated otherwise, $Y$ is finite. When we say \emph{join-semilattice}, we mean every pair $u,v\in Y$ has a least upper bound $u\vee v$.
When we say \emph{monotone cost}, we mean $u\preceq v \Rightarrow c(u)\le c(v)$.
All background on TV-as-testing and the chain/stop-time specialization is in the four prior papers (titles in the abstract).

\vspace{0.25em}\noindent
\textbf{Chain/stop-time $\leftrightarrow$ poset dictionary (quick bridge).}
\begin{itemize}[leftmargin=*]
\item \textbf{Time index / transcript.} ``Time'' in the chain model $\longleftrightarrow$ a transcript $y\in Y$ ordered by add-only feasibility $x\preceq y$.
\item \textbf{Timeout / deadline censoring.} Mapping $x$ to either itself or a terminal symbol $\top$ $\longleftrightarrow$ top-censoring on a poset with maximum element $\top$.
\item \textbf{Prefix sets.} Prefix $\{t'\!\le t\}$ $\longleftrightarrow$ a principal ideal/down-set $\downarrow y=\{u\in Y: u\preceq y\}$.
\item \textbf{Intermediate pooling.} ``Stop early at a common time'' $\longleftrightarrow$ pool incomparable states to a common upper bound $j\succeq x,x'$ (often $j=x\vee x'$).
\item \textbf{Transport lower bounds.} Chain transport/coupling viewpoints $\longleftrightarrow$ join-transport objectives when $Y$ is a join-semilattice.
\item \textbf{Regimes/kinks.} Breakpoints in the chain tradeoff curve $\longleftrightarrow$ changes in the active dual constraints in the compact dual (single $\lambda^\star$).
\end{itemize}

\vspace{0.25em}\noindent
\textbf{Poset archetypes (what to try first).}
\begin{itemize}[leftmargin=*]
\item \textbf{Chain with a terminal top:} you are in the timeout/deadline setting; if top-alignment holds, use the knapsack reduction.
\item \textbf{Diamond/join-semilattice:} incomparable worlds typically benefit from pooling to joins; try the swap-world pooling curve and join optimality first.
\item \textbf{General finite poset:} solve the LP once, then interpret the solution through the compact dual (slope/kink certificates) and sparsity/menu-size bounds.
\item \textbf{Partial observability:} if the adversary only sees features, enforce feature-family TV and compress per-row supports.
\end{itemize}



\vspace{0.25em}\noindent
\textbf{Where the missing pieces live.}
For the testing interpretation of $\TV$ and ``many-hypothesis'' viewpoints, use \emph{Spectral Anonymity and Optimal Distinguishing Bounds for Random-Walk Delegation in (Possibly Directed) Overlays} and \emph{many\_testing\_pub\_tight\_section2and4}.
For knapsack/transport certificates and separation-cut machinery, use \emph{Prefix Capacities and Separation Cuts for Shared-Kernel Termination-Time Hiding}.
For dual-regime reading (slopes/kinks) and refined certificate folklore, use \emph{stop\_time\_unique\_addendum\_v3}.





\vspace{0.25em}\noindent
\textbf{Common pitfalls (read once).}
If $Y$ has no maximum element, top-censoring is not a meaningful baseline. If $Y$ is not a join-semilattice, join pooling is not defined for all pairs (but pooling to \emph{some} common upper bound may still help). If the cost $c$ is not monotone, several ``least upper bound'' arguments fail and you should revert to the LP+dual workflow.

\vspace{0.25em}\noindent
\textbf{From protocol to poset (selection checklist).}
When applying the theory to a real system, the main modeling choice is the poset $(Y,\preceq)$ and the cost $c$.
A reliable workflow is:
\begin{enumerate}[leftmargin=*]
\item \textbf{Choose the transcript space $Y$.} Decide what the adversary actually observes (e.g., termination time buckets, message counts, visible metadata, partial logs). Make each $y\in Y$ a \emph{distinct observable outcome}.
\item \textbf{Define ``add-only'' as a partial order.} Put $x\preceq y$ if transcript $y$ can be produced from intrinsic state $x$ by \emph{only adding} permitted artifacts (padding, delays, dummy events, extra bytes, extra hops). Check transitivity.
\item \textbf{Sanity check joins/tops.} If your system naturally has a ``maximal'' transcript (everything padded to the worst case), that is a $\top$. If your system supports \emph{merging} two outcomes into a least common upper bound, you likely have joins (use pooling/join results).
\item \textbf{Pick a monotone cost $c$.} Encode the operational burden: latency, bandwidth, compute, user-visible friction. Monotonicity ($u\preceq v\Rightarrow c(u)\le c(v)$) is usually appropriate for padding.
\item \textbf{Specify the worlds.} Identify $\pi_0,\pi_1$ (or a family) from the threat model: the two behaviors you want to make hard to distinguish. Swap-world reductions are a first diagnostic; otherwise solve the LP.
\item \textbf{Decide what ``privacy'' measures.} Use transcript TV when the attacker sees the full transcript; use feature-family TV if only coarse features are visible (e.g., counts, bins, summaries).
\end{enumerate}
\noindent
This checklist is intentionally pragmatic: most mistakes we see are mismatched $Y$ (wrong observable) or a non-monotone $c$ (invalidates least-upper-bound arguments).


\section[Aligned top-censoring to fractional knapsack]{Aligned top-censoring $\Rightarrow$ fractional knapsack (recovered special case)}
\noindent\textbf{Role in this addendum.} This recovers the chain/timeout solution as a special case of the poset model and isolates the exact hypothesis (top-alignment) under which it remains optimal.
Assume $(Y,\preceq)$ has a maximum $\top$ and restrict to \emph{top-censoring} kernels supported on $\{x,\top\}$.
Write $s(x):=K(x\mid x)\in[0,1]$ and $\Delta(x):=\pi_0(x)-\pi_1(x)$.

\begin{definition}[Top-alignment]\label{def:topalignment}
The instance is \emph{top-aligned} if $\Delta(x)\ge 0$ for all $x\neq \top$ (equivalently $\Delta(\top)\le 0$).
\end{definition}

\begin{remark}[Is $\top$ an intrinsic state?]
Some variants treat $\top$ as a \emph{real} intrinsic state ($\top\in X$), which makes alignment conditions involving $w(\top)$ meaningful (this matches ``deadline-as-state'' conventions in the chain/stop-time papers).
If in your application $\top$ is \emph{output-only} ($\top\notin X$), then read top-alignment as a condition over $x\in X$ only and ignore any equivalent reformulations that mention $w(\top)$.
\end{remark}


\begin{lemma}[Aligned TV budget in the top-censoring subclass]\label{lem:tvcollapse}
If the instance is top-aligned, then under any top-censoring kernel,
\[
\TV(P_0,P_1)=\sum_{x\neq \top} \Delta(x)\, s(x).
\]
\end{lemma}

\begin{remark}[Intuition for Lemma~\ref{lem:tvcollapse}]
Under top-censoring, the only discrepancy between $P_0$ and $P_1$ that is not ``pinned'' to an intrinsic atom is the mass moved to $\top$.
Top-alignment makes all discrepancies have the same sign on intrinsic atoms, so the $\ell_1$ objective collapses to a single signed sum, i.e.\ one linear budget in the $s(x)$ variables.
\end{remark}



\noindent\textbf{Knapsack value.} For each intrinsic $x\in X\setminus\{\top\}$ define the ``value of not censoring'' 
\[
v(x)\;:=\;\bar\pi(x)\big(c(\top)-c(x)\big).
\]

\begin{theorem}[Aligned knapsack]\label{thm:knapsack}
In the top-aligned regime, optimizing average expected cost over the top-censoring subclass subject to $\TV(P_0,P_1)\le \varepsilon$
is equivalent to fractional knapsack:
\[
\max_{0\le s(x)\le 1}\ \sum_{x\neq \top} v(x)s(x)
\quad \text{s.t.}\quad \sum_{x\neq \top}\Delta(x)s(x)\le \varepsilon.
\]
The greedy policy is optimal when items are sorted by ratio $\rho(x)=v(x)/\Delta(x)$.
\end{theorem}

\noindent\textbf{Greedy/certificates:} see \emph{Prefix Capacities and Separation Cuts for Shared-Kernel Termination-Time Hiding} (and refinements in \emph{stop\_time\_unique\_addendum\_v3}).

\begin{remark}[Why this subsection is here]
This knapsack recovery is included only to pin down the \emph{exact} hypothesis (top-alignment) under which the familiar timeout-style solution is optimal in the poset model.
When top-alignment fails, top-censoring can be strictly suboptimal; Section~2 records the pooling/join phenomenon that replaces it.
\end{remark}


\section{Beyond alignment: pooling on upper bounds and join optimality}
\noindent\textbf{Role in this addendum.} This is the first genuinely poset-native phenomenon: optimal kernels may map both worlds to an intermediate common upper bound (pooling) rather than censoring to $\top$.

\vspace{0.25em}\noindent
\textbf{Swap worlds.} The cleanest setting is when the two priors differ only by swapping mass between two intrinsic states $x$ and $x'$.
In that case, the entire TV constraint is controlled by how different the kernel rows $K(\cdot\mid x)$ and $K(\cdot\mid x')$ are.

\begin{definition}[Swap-world instance]\label{def:swapworld}
Fix two intrinsic states $x,x'\in X$.
A \emph{swap-world} pair $(\pi_0,\pi_1)$ places probability $p$ on $x$ and $1-p$ on $x'$ in world $0$, and swaps them in world $1$:
$\pi_0(x)=p,\ \pi_0(x')=1-p$ and $\pi_1(x)=1-p,\ \pi_1(x')=p$, with $p\in(\tfrac12,1)$ and $\delta:=2p-1$.
\end{definition}

\begin{lemma}[TV reduces to row-wise TV]\label{lem:tvrows}
In a swap-world instance on $\{x,x'\}$, for any poset-feasible kernel $K$,
\[
\TV(P_0,P_1)=\Delta\cdot \TV\big(K(\cdot\mid x),\,K(\cdot\mid x')\big).
\]
\end{lemma}

\begin{remark}[Intuition for Lemma~\ref{lem:tvrows}]
In a swap world, all prior differences concentrate on $x$ versus $x'$.
After applying the same kernel, the induced transcript difference is just that signed mass times the difference of the two kernel rows, so $\TV(P_0,P_1)$ becomes proportional to $\|K(\cdot\mid x)-K(\cdot\mid x')\|_1$.
\end{remark}


\begin{theorem}[Optimal pooling within $\{x,x',j\}$]\label{thm:pooling}
For any $\varepsilon\in[0,\Delta]$, the choice $s=\varepsilon/\Delta$ minimizes average expected cost among kernels supported on $\{x,x',j\}$, and
\[
\TV(P_0,P_1)=\Delta s,\qquad
\mathrm{OPT}(\varepsilon)=c(j)-\frac{\varepsilon}{\Delta}\Big(c(j)-\tfrac12(c(x)+c(x'))\Big).
\]
In particular, perfect privacy ($\varepsilon=0$) is achieved at cost $c(j)$, potentially much smaller than padding to a global top.
\end{theorem}

\begin{theorem}[Pooling on the join is globally optimal for swap worlds]\label{thm:joinglobal}
Assume $(Y,\preceq)$ is a finite \emph{join-semilattice} and the cost $c$ is monotone: $u\preceq v\Rightarrow c(u)\le c(v)$.
Let $x,x'\in X$ be incomparable and let $j=x\vee x'$ be their join (least common upper bound).
In a swap-world instance supported on $\{x,x'\}$, for any privacy budget $\varepsilon\in[0,\Delta]$ the minimum average expected cost under
$\TV(P_0,P_1)\le \varepsilon$ equals
\[
\mathrm{OPT}(\varepsilon)=c(j)-\frac{\varepsilon}{\Delta}\Big(c(j)-\tfrac12(c(x)+c(x'))\Big),
\]
and is achieved by the pooling kernel
\[
K(j\mid x)=K(j\mid x')=1-\tfrac{\varepsilon}{\Delta},\qquad K(x\mid x)=K(x'\mid x')=\tfrac{\varepsilon}{\Delta}.
\]
In particular, allowing arbitrary feasible outputs beyond $\{x,x',j\}$ cannot improve the optimum.
\end{theorem}

\begin{example}[Micro-example: diamond pooling beats top-censoring]
Let $Y=\{x,x',j,\top\}$ with $x\not\preceq x'$, $x'\not\preceq x$, $x\preceq j\preceq \top$, and $x'\preceq j\preceq \top$.
Assume monotone costs $c(x)=c(x')=0$, $c(j)=1$, $c(\top)=M$ with $M\gg 1$.
Consider swap worlds supported on $\{x,x'\}$: $\pi_0(x)=\tfrac12+\delta$, $\pi_0(x')=\tfrac12-\delta$, and $\pi_1$ swaps the masses.
\begin{itemize}[leftmargin=*]
\item \textbf{Top-censoring.} Any top-censoring kernel that reduces $\TV(P_0,P_1)$ must send a nontrivial fraction of both $x$ and $x'$ to $\top$, incurring cost on the order of $M$.
\item \textbf{Join pooling.} Pooling both states to $j$ with probability $1-s$ makes the induced TV budget affine in $s$ (Theorem~\ref{thm:pooling}), so one can choose $s$ to hit a target $\varepsilon$, while paying only $c(j)=1$ when pooling occurs.
When $M$ is large, this yields a strict improvement over top-censoring at the same $\varepsilon$.
\end{itemize}
This is the smallest poset instance where ``censor-to-top'' is dominated by pooling at an intermediate upper bound.
\end{example}


\section[Perfect transcript privacy (epsilon=0)]{Perfect transcript privacy ($\varepsilon=0$): OT relaxation and tightness}

\begin{remark}[A three-point mental model for pooling]
In a diamond-shaped subposet with two incomparable points $x,x'$ and an upper bound $j\succeq x,x'$, pooling means:
``leave the input alone with probability $s$, otherwise output $j$.'' 
Theorems in this section state that (a) the privacy cost is affine in $s$ for swap worlds, and (b) on join-semilattices with monotone cost, choosing $j=x\vee x'$ is globally optimal.
\end{remark}

We record the three compact theorems needed later (dual, join-transport relaxation, and a sharp implementability criterion).

\begin{align}
\min_{K\ge 0}\quad &\sum_{x\in X}\bar\pi(x)\sum_{y\succeq x} c(y)\,K(y\mid x) \label{eq:pp_primal}\\
\text{s.t.}\quad &\sum_{x\in X}(\pi_0(x)-\pi_1(x))\,K(y\mid x)=0\quad \forall y\in Y, \label{eq:pp_balance}\\
&\sum_{y\succeq x}K(y\mid x)=1\quad \forall x\in X. \nonumber
\end{align}

\begin{remark}[Relaxation gaps can be real]
The join-transport relaxation is a lower bound on the world-oblivious optimum. It can be tight, but it can also be strictly optimistic when likelihood-ratio invariance fails.
The full version contains explicit gap constructions (including chain examples where the gap can be unbounded).
\end{remark}
\begin{theorem}[Dual of world-oblivious perfect privacy]\label{thm:pp_dual}
The dual of~\eqref{eq:pp_primal}--\eqref{eq:pp_balance} is
\begin{align}
\max_{\alpha,\beta}\quad & \sum_{x\in X}\alpha(x)\label{eq:pp_dual_obj}\\
\text{s.t.}\quad & \alpha(x)+d(x)\,\beta(y)\le \bar\pi(x)\,c(y)\qquad \forall (x,y)\ \text{with}\ x\preceq y.\label{eq:pp_dual_constr}
\end{align}
Strong duality holds. In particular, any dual-feasible pair \((\alpha,\beta)\) yields the lower bound
\[
\mathrm{Opt}_{\varepsilon=0}\ \ge\ \sum_{x\in X}\alpha(x).
\]
\end{theorem}

\begin{theorem}[Join-transport value for the coupling relaxation]\label{thm:jointransport_relax}
Assume \((Y,\preceq)\) is a finite join-semilattice and \(c\) is monotone.
In the two-kernel relaxation described above, the minimum cost at \(\varepsilon=0\) equals
\[
\inf_{\Gamma\in \Pi(\pi_0,\pi_1)} \E_{(X,X')\sim \Gamma}\big[c(X\vee X')\big],
\]
where \(\Pi(\pi_0,\pi_1)\) is the set of couplings of \(\pi_0\) and \(\pi_1\).
\end{theorem}

\begin{theorem}[Tightness criterion via posterior likelihood ratios]\label{thm:tightness}
Fix \(\pi_0,\pi_1\) on \(X\) and a poset-feasible transcript space \(Y\).
A distribution \(Q\) on \(Y\) is achievable with \(\varepsilon=0\) by a world-oblivious kernel \(K\) (i.e.\ \(P_0=P_1=Q\))
if and only if there exist posteriors \(\eta_0(\cdot\mid y),\eta_1(\cdot\mid y)\) such that:
\begin{enumerate}
\item (Marginals) \(\sum_y Q(y)\eta_b(x\mid y)=\pi_b(x)\) for \(b\in\{0,1\}\);
\item (Poset support) \(\eta_b(x\mid y)=0\) unless \(x\preceq y\);
\item (Likelihood-ratio invariance) \(\eta_0(x\mid y)\pi_1(x)=\eta_1(x\mid y)\pi_0(x)\) for all \(x,y\) with \(\pi_0(x)\pi_1(x)>0\).
\end{enumerate}
Moreover, when these conditions hold, the corresponding world-oblivious kernel is given by
\[
K(y\mid x)=\frac{Q(y)\eta_0(x\mid y)}{\pi_0(x)}=\frac{Q(y)\eta_1(x\mid y)}{\pi_1(x)}
\quad\text{for all }x\text{ with }\pi_0(x)\pi_1(x)>0,
\]
with arbitrary choices on \(x\) outside the support of one prior.
\end{theorem}

\begin{remark}[How to check likelihood-ratio invariance in practice]
Compute (or simplify) the posterior likelihood ratio implied by the relaxed construction at each transcript $y$.
The tightness theorem says the relaxation is implementable by a \emph{single} world-oblivious kernel exactly when these Bayes-factor ratios can be made consistent across worlds (informally: ``$y$ carries the same evidence in both constructions'').
\end{remark}


\section[General epsilon: compact dual and sparsity]{General $\varepsilon$: compact dual and sparsity/menu-size consequences}
\noindent\textbf{Role in this addendum.} Let $F(\varepsilon)$ denote the optimal expected cost under $\TV(P_0,P_1)\le \varepsilon$. By LP duality, $F$ is convex and piecewise-linear; the compact dual makes the active ``regime'' visible via a single multiplier $\lambda^\star$.
\vspace{0.25em}
\noindent\textbf{Slope certificate (display).} For any $\varepsilon_0$, if $\lambda^\star$ is optimal at $\varepsilon_0$, then $F(\varepsilon)\ge F(\varepsilon_0)-2\lambda^\star(\varepsilon-\varepsilon_0)$ for all $\varepsilon$.



\subsection{Transcript-TV primal LP}\label{sec:lp_global}
Fix world priors $\pi_0,\pi_1$ on intrinsic states $X$, let $d(x)=\pi_0(x)-\pi_1(x)$ and $\bar\pi(x)=(\pi_0(x)+\pi_1(x))/2$.
A world-oblivious padding kernel $K(\cdot\mid x)$ supported on monotone transcripts $y\succeq x$ induces output laws
$P_s(y)=\sum_x \pi_s(x)K(y\mid x)$.
The global design problem with transcript total-variation budget $\TV(P_0,P_1)\le \varepsilon$ is the LP
\begin{align}
\min_{K\ge 0}\quad & \sum_{x\in X}\sum_{y\succeq x} \bar\pi(x)\,c(y)\,K(y\mid x)\\
\text{s.t.}\quad & \sum_{y\succeq x} K(y\mid x)=1\qquad \forall x\in X,\\
& \frac12\sum_{y\in Y}\left|\sum_{x\in X} d(x)\,K(y\mid x)\right|\le \varepsilon.
\end{align}
This is the primal whose compact dual is stated in Theorem~\ref{thm:dual_general_eps}.

\begin{theorem}[Dual for \(\TV(P_0,P_1)\le \varepsilon\)]\label{thm:dual_general_eps}
The optimal value \(F(\varepsilon)\) of the world-oblivious transcript-TV design problem equals
\begin{align}
F(\varepsilon)=\sup_{\alpha,\beta,\lambda}\quad & \sum_{x\in X}\alpha(x)-2\varepsilon\,\lambda \label{eq:dual_eps_obj}\\
\text{s.t.}\quad & \alpha(x)+d(x)\,\beta(y)\le \bar\pi(x)\,c(y)\qquad \forall (x,y)\ \text{with}\ x\preceq y, \label{eq:dual_eps_constr}\\
& |\beta(y)|\le \lambda\qquad \forall y\in Y,\qquad \lambda\ge 0. \label{eq:dual_eps_box}
\end{align}
Strong duality holds. In particular, any dual-feasible \((\alpha,\beta,\lambda)\) certifies the lower bound \(F(\varepsilon)\ge \sum_x \alpha(x)-2\varepsilon\lambda\).
\end{theorem}

\begin{remark}[Workflow: solve once, then certify]
Solve the primal LP at a target $\varepsilon$, extract an optimal dual solution and the multiplier $\lambda^\star$, and read off the active constraints to identify which transcripts/rows govern privacy at that budget.
Scanning a small grid of $\varepsilon$ values reveals kink locations (regime changes) cheaply.
\end{remark}



\noindent\textbf{Mixture viewpoint (shared randomness).}
We view a kernel as a \emph{convex combination of deterministic padding maps} (a distribution over monotone maps $f:X\to Y$ with $x\preceq f(x)$),
so the same shared randomness selects a single map and applies it to all $x$; this is the representation under which basic-feasible-solution / Carathéodory support bounds yield small menus.

\begin{theorem}[Menu-size bound: few deterministic paddings suffice]\label{thm:menu_size_det}
In the transcript-TV LP (Section~\ref{sec:lp_global}), there exists an optimal kernel that can be written as a mixture of at most \(2|Y|+1\) deterministic paddings:
\[
K^\star=\sum_{i=1}^m w_i K_{f_i},\qquad m\le 2|Y|+2,\quad w_i\ge 0,\quad \sum_i w_i=1.
\]
Equivalently, optimal padding can be implemented by sampling one of \(m\) deterministic padding maps and applying it to \(x\).
\end{theorem}

\section{Feature-family TV (optional): per-row support compression}
\noindent\textbf{Drop-in module.} Keep this only if you enforce privacy on a finite observable feature family; otherwise delete this section with no dependencies.


\begin{remark}[How to read regimes from the dual]
In the compact dual (Theorem~\ref{thm:dual_general_eps}), the optimal multiplier $\lambda^\star$ is the slope of the supporting hyperplane of the tradeoff curve at $\varepsilon$.
As $\varepsilon$ varies, kinks occur exactly when a different subset of dual constraints becomes active.
This ``regime certificate'' viewpoint is developed in detail in \emph{stop\_time\_unique\_addendum\_v3}; we only record the dual statement here.
\end{remark}

\begin{theorem}[Per-row support compression]\label{thm:caratheodory}
Under feature-family TV constraints and the average-cost objective, there exists an optimal kernel $K^\star$ such that for every intrinsic state $x$,
\[
|\mathrm{supp}(K^\star(\cdot\mid x))|\le d+2,\qquad d=\sum_{f\in\mathcal{F}} k_f.
\]
\end{theorem}

\section*{Checklist (one box)}
\begin{framed}
\noindent\textbf{Use this addendum as follows.}
\begin{enumerate}[leftmargin=*]
\item \textbf{Identify the geometry.} Specify $(Y,\preceq)$ (what counts as ``add-only'') and a monotone cost $c$ on transcripts. Decide whether $Y$ has a maximum $\top$ and/or joins $x\vee x'$.
\item If you are in the top-censoring/timeout subclass: verify top-alignment; apply knapsack; use greedy + dual certificates from \emph{Prefix Capacities and Separation Cuts for Shared-Kernel Termination-Time Hiding}.
\item For swap worlds on join-semilattices: pool to the join (global optimality).
\item For $\varepsilon=0$: compute OT relaxation value; check likelihood-ratio invariance for implementability.
\item For general $\varepsilon$: solve the LP once; read $\lambda^\star$ as slope/kink certificate; use sparsity/menu-size bounds to simplify implementation.
\end{enumerate}
\end{framed}
\vspace{0.25em}\noindent
\textbf{How the threads connect (one paragraph).}
Top-aligned censoring collapses the privacy constraint to a single ``prefix'' budget and recovers the knapsack structure familiar from the chain/stop-time setting.
When alignment fails, the right primitive is row-wise coupling: swap worlds reduce the constraint to a distance between two kernel rows, which makes intermediate pooling natural and, on join-semilattices, globally optimal at the join.
At $\varepsilon=0$, allowing separate kernels per world yields an OT-style lower bound (join-transport); the likelihood-ratio invariance condition pins down exactly when a \emph{single} world-oblivious kernel achieves it.
For general $\varepsilon$, the compact dual packages all of this into a single multiplier $\lambda^\star$ whose value certifies slopes/kinks and enables sparsity-based implementation.




\section[Practical strategies]{Practical strategies and what this points toward (with light speculation)}
\noindent
This addendum is deliberately theorem-forward, but a few practical design patterns fall out of the poset view.

\subsection*{Design patterns suggested by the poset model}
\begin{itemize}[leftmargin=*]
\item \textbf{Prefer intermediate pooling states over ``pad to worst case.''}
If your transcript space has meaningful intermediate upper bounds (joins or near-joins) with much lower cost than $\top$, then top-censoring can be dominated.
A concrete engineering interpretation is: add an intermediate \emph{aggregation/merge} outcome (a shared ``summary'' transcript) so that incomparable behaviors can be pooled cheaply.
\item \textbf{Make joins real by adding mergeable abstractions.}
When joins do not exist, you can sometimes \emph{create} them by adding a coarser, mergeable representation (e.g., bucketed times, capped counters, quantized sizes, or standardized dummy sequences) that both worlds can map to.
This is a principled version of ``design a common padded form that many behaviors can share.''
\item \textbf{Use the dual as an operational dial.}
The compact dual produces a single multiplier $\lambda^\star$ at a chosen $\varepsilon$.
Operationally, $\lambda^\star$ is the marginal cost of tightening privacy; kinks indicate regime changes (which rows/transcripts matter).
This supports a workflow like: pick an $\varepsilon$ at a kink, where additional privacy starts becoming expensive, and treat $\lambda^\star$ as a tuning knob.
\item \textbf{Exploit small menus for implementability.}
Sparsity/menu-size theorems say optimal kernels can be implemented as mixtures of few deterministic padding maps.
That suggests caching and simple policy tables rather than heavy online optimization, and it reduces the surface area for implementation bugs.
\item \textbf{Model what the attacker \emph{actually} sees.}
If only coarse features are observable, enforce feature-family TV instead of full transcript TV and compress per-row supports.
This can justify cheaper defenses that would be ``overkill'' under full-transcript observability.
\end{itemize}

\subsection*{How this might apply ``in the world'' (examples)}
\begin{itemize}[leftmargin=*]
\item \textbf{Network protocols / timing side channels.} Replace ``always delay to max'' with a small set of shared delay profiles (intermediate pooling states), chosen by a sparse mixture, calibrated to a TV budget.
\item \textbf{Logging/telemetry.} Introduce mergeable summary events (joins) so that multiple user behaviors map to a common high-level transcript while adding only allowable extra fields/events.
\item \textbf{User-facing privacy knobs.} Use the dual slope $\lambda^\star$ to report a meaningful ``cost per unit privacy'' to product owners, and to choose stable operating points away from unstable regimes.
\end{itemize}

\subsection*{Light speculation: shaping the poset is part of the defense}
A recurring lesson is that privacy is not only about choosing $K$ \emph{given} $Y$; it is also about designing the transcript space itself.
If you can redesign instrumentation so that (i) add-only padding is natural, (ii) monotone costs reflect reality, and (iii) intermediate mergeable transcripts exist, then the optimal defenses tend to become both cheaper and simpler.
This suggests a co-design loop: pick a candidate $(Y,\preceq,c)$, solve the LP on representative world pairs, inspect the dual to see which transcripts drive the privacy constraint, and then adjust what is logged/observable to create cheaper pooling opportunities.


\subsection*{Deployment cautions (read once)}
\begin{framed}
\noindent\textbf{Cautions.} These are not theorems---just failure modes we have repeatedly seen in real deployments of padding/obfuscation mechanisms.
\begin{itemize}[leftmargin=*]
\item \textbf{Attacker model mismatch.} If the attacker observes richer data than your modeled $Y$ (e.g., fine-grained timing instead of bins), guarantees can evaporate. Prefer modeling the \emph{richest} plausible observable transcript, then downgrade to feature-family TV only when justified.
\item \textbf{Composability.} Padding decisions made per-session can leak when aggregated across sessions (linkage, averaging, correlation). If composition matters, validate the effective privacy budget under repeated use.
\item \textbf{Side-channel ``leaks around'' the poset.} The model assumes add-only feasibility. Real systems may introduce non-monotone artifacts (timeouts, retries, caching, queueing effects) that violate $x\preceq y$ or cost monotonicity; monitor and patch these first.
\item \textbf{Randomness quality and determinism.} Sparse mixtures are great for implementability, but beware deterministic tie-breaking, shared seeds, or low-entropy PRNGs that let an attacker condition on hidden structure.
\item \textbf{Regime instability near kinks.} Operating exactly at a kink can be sensitive to drift (changes in priors, workloads, or costs). Use the dual slope $\lambda^\star$ to pick stable operating points with margin.
\item \textbf{Overfitting to a few world pairs.} Optimizing for one $(\pi_0,\pi_1)$ can shift leakage to other plausible alternatives. Stress-test against a \emph{family} of adversary hypotheses (see the many-testing viewpoint in the prior papers).
\item \textbf{Human factors.} If padding produces user-visible delays or artifacts, users may adapt behavior and change $\pi_b$; incorporate feedback loops when estimating priors.
\end{itemize}
\end{framed}

\noindent\textbf{Citation practice.}
When reusing knapsack/greedy/transport proofs or dual-certificate arguments, cite \emph{Prefix Capacities and Separation Cuts for Shared-Kernel Termination-Time Hiding}.
When reusing the TV-as-testing and hypothesis-testing framing, cite \emph{Spectral Anonymity and Optimal Distinguishing Bounds for Random-Walk Delegation in (Possibly Directed) Overlays} and/or \emph{many\_testing\_pub\_tight\_section2and4}.
When reusing kink/slope regime-reading or refined dual folklore, cite \emph{stop\_time\_unique\_addendum\_v3}.
\vspace{0.25em}

\section*{References (titles only)}
{\small (Assumed available to the reader; this addendum is intentionally not self-contained.)}
\begin{itemize}[leftmargin=*]
\item \emph{Spectral Anonymity and Optimal Distinguishing Bounds for Random-Walk Delegation in (Possibly Directed) Overlays}.
\item \emph{Prefix Capacities and Separation Cuts for Shared-Kernel Termination-Time Hiding}.
\item \emph{stop\_time\_unique\_addendum\_v3}.
\item \emph{many\_testing\_pub\_tight\_section2and4}.
\end{itemize}

\end{document}
